% ============================================================
%   Exercice Simplexe - Sujet 1 : Transport de poissons
%   Fichier LaTeX autonome, compilable directement par :
%       pdflatex simplexe.tex
% ============================================================
\documentclass[12pt,a4paper]{article}

% --- Encodage et langue ---
\usepackage[utf8]{inputenc}
\usepackage[T1]{fontenc}
\usepackage{lmodern}
\usepackage[french]{babel}

% --- Mathematiques ---
\usepackage{amsmath}
\usepackage{amssymb}
\usepackage{amsfonts}
\usepackage{mathtools}

% --- Mise en page ---
\usepackage{geometry}
\geometry{margin=2.5cm}
\usepackage{parskip}

% ============================================================
%   COMMANDE \Eq : signe egal renforce (visible dans tout lecteur
%   PDF, y compris VS Code PDF Viewer). Police texte + Large + bold.
% ============================================================
\newcommand{\Eq}{\mathrel{\text{\Large\bfseries =}}}

% ============================================================
%   COMMANDES \Leq et \Geq : inegalites <= et >= en version
%   inclinee (plus visible et bien distincte de < et >).
% ============================================================
\newcommand{\Leq}{\leqslant}
\newcommand{\Geq}{\geqslant}

% --- Equation encadree avec \Eq ---
\newcommand{\BoxedEq}[1]{%
  \begin{equation*}
    \boxed{\,\displaystyle #1\,}
  \end{equation*}%
}

\begin{document}

\section*{Exercice Simplexe : Transport de poissons}

\subsection*{Contexte}
Une entreprise souhaite organiser le transport du poisson frais entre Nouadhibou et Nouakchott. On considere deux destinations principales : le marche principal et les restaurants/supermarches.

\subsection*{Variables de decision}
\[
\begin{aligned}
X &\Eq \text{quantite de poisson transportee vers le marche principal}, \\
Y &\Eq \text{quantite de poisson transportee vers les restaurants et supermarches}.
\end{aligned}
\]

\subsection*{Fonction objectif}
L'objectif est de maximiser le benefice total :

\BoxedEq{Z_{\mathrm{max}} \Eq 40\,X + 50\,Y}

\subsection*{Contraintes}
\[
\left\{
\begin{aligned}
X + Y    &\Leq 100 \\
2X + 3Y  &\Leq 240 \\
X + 2Y   &\Leq 160 \\
X, Y     &\Geq 0
\end{aligned}
\right.
\]

\subsection*{Forme standard}
On ajoute les variables d'ecart \(S_1\), \(S_2\) et \(S_3\) :
\[
\begin{aligned}
X + Y + S_1            &\Eq 100 \\
2X + 3Y + S_2          &\Eq 240 \\
X + 2Y + S_3           &\Eq 160 \\
Z - 40X - 50Y          &\Eq 0
\end{aligned}
\]

\subsection*{Resolution}
Les sommets realisables sont :
\[
(0,0), \quad (100,0), \quad (0,80), \quad (60,40)
\]

Calcul de la fonction objectif a chaque sommet :
\[
\begin{aligned}
Z(0,0)     &\Eq 0 \\
Z(100,0)   &\Eq 40(100) + 50(0)  \Eq 4000 \\
Z(0,80)    &\Eq 40(0) + 50(80)   \Eq 4000 \\
Z(60,40)   &\Eq 40(60) + 50(40)  \Eq 4400
\end{aligned}
\]

\subsection*{Solution optimale}

\BoxedEq{X \Eq 60 \quad\text{et}\quad Y \Eq 40}

\BoxedEq{Z_{\mathrm{max}} \Eq 4400}

\subsection*{Interpretation}
L'entreprise doit transporter 60 unites de poisson vers le marche principal et 40 unites vers les restaurants et supermarches. Le benefice maximal obtenu est de 4400 unites.

\end{document}
